\chapter{Transition Functions}

\section{Motivation}

\begin{remark}
Transitions are the decisive element type for curvature-driven alignment
modelling. While fixed elements describe geometric states of constant
curvature, transitions describe the controlled change between such states.
\end{remark}

\begin{remark}
Their importance is both mathematical and engineering-related.
Mathematically, they define how curvature evolves. In engineering
contexts, they control smoothness, feasibility, comfort, and
constructability.
\end{remark}

\section{Definition by Curvature Evolution}

\begin{definition}[Transition]
A transition is an alignment element whose curvature varies over its
parameter interval.
\end{definition}

\begin{definition}[Transition curvature law]
Let $L>0$. A transition curvature law is a function
\[
\kappaf : [0,L] \to \R
\]
with prescribed boundary values
\[
\kappaf(0)=\kappa_0,
\qquad
\kappaf(L)=\kappa_1.
\]
\end{definition}

\begin{remark}
The transition geometry is determined by the curvature law together with
the reconstruction principles introduced in the axiomatic foundation.
\end{remark}

\section{Normalized Representation}

\begin{definition}[Normalized parameter]
A normalized local parameter is a variable
\[
u \in [0,1].
\]
\end{definition}

\begin{definition}[Normalized transition law]
A normalized transition law is a function
\[
\widehat{\kappaf} : [0,1] \to \R
\]
with prescribed boundary conditions, typically
\[
\widehat{\kappaf}(0)=0,
\qquad
\widehat{\kappaf}(1)=1.
\]
\end{definition}

\begin{remark}
Normalization separates the abstract shape of the transition from its
physical scale. This makes transition families reusable and comparable.
\end{remark}

\section{Scaling to Physical Transitions}

\begin{proposition}
Let $\widehat{\kappaf}$ be a normalized transition law on $[0,1]$.
For given $L>0$ and curvature values $\kappa_0,\kappa_1 \in \R$, the
corresponding physical transition law is given by
\[
\kappaf(s)
=
\kappa_0 + (\kappa_1-\kappa_0)\,
\widehat{\kappaf}\!\left(\frac{s}{L}\right),
\qquad s\in[0,L].
\]
\end{proposition}

\begin{remark}
This formula separates:
\begin{itemize}
  \item shape, encoded by $\widehat{\kappaf}$,
  \item length scale, encoded by $L$,
  \item and curvature scale, encoded by $\kappa_1-\kappa_0$.
\end{itemize}
\end{remark}

\section{Regularity and Smoothness}

\begin{definition}[Regularity class]
A transition law belongs to class $C^k$ if its curvature function is
$k$-times continuously differentiable.
\end{definition}

\begin{remark}
In practice, the relevance of a transition law is not determined only by
$\kappaf$ itself, but also by the behaviour of its derivatives such as
$\kappaf'$ and $\kappaf''$.
\end{remark}

\OPEN{Formal role of $\kappaf'$ and $\kappaf''$ in engineering semantics, comfort rules, and construction logic.}

\section{Transition Families}

\begin{definition}[Transition family]
A transition family is a class of normalized transition laws sharing a
common structural principle.
\end{definition}

\begin{remark}
Different families differ not only in formulae, but in regularity,
boundary behaviour, implementation complexity, and engineering
interpretation.
\end{remark}

\subsection{Classical Families}

\TODO{Clothoid as the classical case of linear curvature evolution.}

\TODO{Polynomial transition families.}

\TODO{Trigonometric transition families.}

\subsection{Extended and Hybrid Families}

\TODO{Vienna-type transition families.}

\TODO{Lookup-based transition families.}

\TODO{Hybrid or composite transition families.}

\RESEARCH{Classification of admissible transition families under geometric and engineering criteria.}

\section{Decomposition Schemes}

\begin{remark}
A transition need not always be treated as a single indivisible object.
It may be useful to represent it as a composition of sub-elements.
\end{remark}

\TODO{Formalize the halfWave -- clothoid -- halfWave decomposition currently guiding the practical development.}

\OPEN{Clarify whether such decompositions belong to the mathematical model itself or only to a constructive implementation layer.}

\section{Operator View}

\begin{remark}
From a structural perspective, a transition may be interpreted as an
operator that transforms one curvature state into another through a
prescribed evolution law.
\end{remark}

\RESEARCH{Develop a formal operator view of transition composition.}

\section{Representation in Practical Systems}

\begin{remark}
In computational settings, transitions may be represented in different
ways, including:
\begin{itemize}
  \item closed analytic laws,
  \item polynomial approximations,
  \item lookup tables,
  \item mixed or compiled representations.
\end{itemize}
\end{remark}

\TODO{Formalize the registry-based transition concept.}

\OPEN{Trade-off between analytic precision, numerical robustness, and evaluation speed.}

\section{Relation to the Sparse Alignment Model}

\begin{remark}
Within the sparse alignment model, transitions are one of the two
fundamental element classes besides fixed elements.
\end{remark}

\begin{remark}
They provide the mechanism for smooth geometric change and therefore
govern the qualitative behaviour of the alignment as a whole.
\end{remark}

\section{Summary}

\begin{summary}
A transition is a curvature-driven alignment element with non-constant
curvature. Its normalized representation separates shape from scale and
thereby enables reusable families and implementation-independent
comparison.

Transitions form the central mechanism by which smoothness and controlled
change enter the sparse alignment model.
\end{summary}
